\documentclass[../../../main.tex]{subfiles}

\begin{document}

\section{Quick refresher on spin}
\label{sec:spin-refresher}

Spin can be described with the operators $\hat S_x,\hat S_y,\hat S_z$ that obey the following
commutation relations:
\begin{equation}
\boxed{\,[\hat S_j,\hat S_k]=i\hbar\,\epsilon_{jk\ell}\hat S_\ell\,}\label{eq:spin-comm}
\qquad(\hbar=1),
\end{equation}
with $i,j,k=x,y,z=1,2,3$ and using the Einstein summation convention.
One can define raising and lowering operators as
\begin{equation}
\boxed{\,\hat S_\pm=\hat S_x\pm i\hat S_y\,}\qquad\text{with}\qquad
\boxed{\,[S_-,S_+]=-2S_z\,}.
\end{equation}

\begin{proof}
\begin{align}
[S_-,S_+]&=[S_x-iS_y,\ S_x+iS_y]\notag\\
&=
[S_x,S_x] +i[S_x,S_y]-i[S_y,S_x]+[S_y,S_y]\notag\\
&=2i[S_x,S_y]
\overset{\eqref{eq:spin-comm}}{=}2i\cdot i\,\epsilon_{12\ell}S_\ell
=-2\,\epsilon_{12\ell}S_\ell
=-2S_z,
\end{align}
where the first step used the bilinearity of the commutator and the last step uses $\epsilon_{jk\ell}=0$ if two indices are equal, leaving
$\epsilon_{123}S_3=S_z$.
\end{proof}

For spin-$\tfrac12$ particles, one can write
\begin{equation}
S_j=\frac{\hbar}{2}\,\sigma_j\qquad\text{with }\sigma_j\text{ the Pauli matrices, }j=x,y,z,
\end{equation}
and often again $\hbar=1$.

\end{document}
